\scanpage{28}
\mspage{22}

Notons d'abord, pour nous y reconnaître,
que les relations b$'$) du $2^\circ$, et les
relations transposées (compte tenu que $\mathrm{Cof}_W$
et $W^{\mathrm{univ}}$ sont closes à g., par $1^\circ$), \deleted{\unsure{donnent}}
savoir
\[
 \left\{\begin{aligned}
 (\mathrm{Cof}_W)_*&\subset W_{\mathrm{univ}},&
 (W^{\mathrm{univ}})_*&\subset\mathrm{Fib}_W\\
 (W_{\mathrm{univ}})^{\unsure{\alpha|*}}&\subset\mathrm{Cof}_W,&
 (\mathrm{Fib}_W)^{\alpha}&\subset W^{\mathrm{univ}}
 \end{aligned}\right.
\]
donnent\textsuperscript{$*$} les deux diagrammes \inserted{carrés} d'inclusions
% slanted
\margin{compte tenu de $W^{\mathrm{univ}}\subset\mathrm{Cof}_W$ $W_{\mathrm{univ}}\subset\mathrm{Fib}_W$ \unsure{$\}$} et des inclusions transposées}

\diagram{\begin{tikzcd}[ampersand replacement=\&,row sep=small,column sep=small]
 \& W^{\mathrm{univ}}\arrow[dr,hook] \& \\
 (\mathrm I)\qquad(\mathrm{Fib}_W)^*\arrow[ur,hook]\arrow[dr,hook] \& \& \mathrm{Cof}_W \\
 \& (W_{\mathrm{univ}})^{\unsure{\wedge|*}}\arrow[ur,hook] \&
\end{tikzcd}}

et
% slanted
\margin{\uline{NB} On a ces carrés d'inclusion, chaque fois qu'on a un syst. de Quillen $(W,C,F,TC,TF)$ \uline{clos}}

\diagram{\begin{tikzcd}[ampersand replacement=\&,row sep=small,column sep=small]
 \& (W^{\mathrm{univ}})_*\arrow[dl,hook,very thick] \& \\
 (\mathrm{II})\qquad\mathrm{Fib}_W \& \& (\mathrm{Cof}_W)_*\arrow[ul,hook]\arrow[dl,hook,very thick] \\
 \& W_{\mathrm{univ}}\arrow[ul,hook,"\alpha"'] \& 
\end{tikzcd}},

dont le premier est transposé (par $\unsure{\Phi|\underline{\varphi}|\underline{\mathcal{Q}}|\underline{\mathcal{D}}}\unsure{-|\to}\Phi^*$)
du second, et où les 8 \unsure{'}ensembles de
flèches envisagés, à l'exception de $\mathrm{Fib}_W$ et
$W_{\mathrm{univ}}$ (reliés par $\alpha$ dans le carré $\mathrm{II}$)
sont clos. (\deleted{\unsure{les} \unsure{de} \unsure{parties}} \inserted{les quatre} ensembles de
$\mathrm I$ clos à g., et les deux ens. de $\mathrm{II}$
non \unsure{intéressés} par $\alpha$, clos à dr.).
% archive sheet number: 27
